\documentclass[11pt]{article}
\usepackage[a4paper,margin=2.1cm]{geometry}
\usepackage{amsmath,amssymb}
\usepackage[T1]{fontenc}
\usepackage{xcolor}
\usepackage{enumitem}
\newcommand{\rev}[1]{#1}% captions of the manuscript contain \rev; they are read from its .aux file
\usepackage{xr-hyper}
\externaldocument{manuscript_r1_clean}
\usepackage[hidelinks]{hyperref}
% generated by make_tex_r1.py -- do not edit
\makeatletter
\newcommand{\N}[1]{\@ifundefined{num@#1}{\textbf{??[#1]}}{\csname num@#1\endcsname}}
\makeatother
\expandafter\def\csname num@cheb.compat.0.01.N120.dt5e-05.l2\endcsname{$5.37\times10^{-6}$}
\expandafter\def\csname num@cheb.compat.0.01.N40.dt0.002.l2\endcsname{$2.21\times10^{-4}$}
\expandafter\def\csname num@cheb.compat.0.01.N60.dt0.001.l2\endcsname{$1.07\times10^{-4}$}
\expandafter\def\csname num@cheb.compat.0.01.N80.dt0.0001.l2\endcsname{$1.07\times10^{-5}$}
\expandafter\def\csname num@cheb.compat.0.01.N80.dt0.0005.l2\endcsname{$5.37\times10^{-5}$}
\expandafter\def\csname num@cheb.compat.0.1.N120.dt5e-05.l2\endcsname{$5.10\times10^{-6}$}
\expandafter\def\csname num@cheb.compat.0.1.N40.dt0.002.l2\endcsname{$2.04\times10^{-4}$}
\expandafter\def\csname num@cheb.compat.0.1.N60.dt0.001.l2\endcsname{$1.02\times10^{-4}$}
\expandafter\def\csname num@cheb.compat.0.1.N80.dt0.0001.l2\endcsname{$1.02\times10^{-5}$}
\expandafter\def\csname num@cheb.compat.0.1.N80.dt0.0005.l2\endcsname{$5.10\times10^{-5}$}
\expandafter\def\csname num@cheb.orig.0.01.N120.dt5e-05.l2\endcsname{$4.22\times10^{-6}$}
\expandafter\def\csname num@cheb.orig.0.01.N40.dt0.002.l2\endcsname{$1.69\times10^{-4}$}
\expandafter\def\csname num@cheb.orig.0.01.N60.dt0.001.l2\endcsname{$8.43\times10^{-5}$}
\expandafter\def\csname num@cheb.orig.0.01.N80.dt0.0001.l2\endcsname{$8.43\times10^{-6}$}
\expandafter\def\csname num@cheb.orig.0.01.N80.dt0.0005.l2\endcsname{$4.22\times10^{-5}$}
\expandafter\def\csname num@cheb.orig.0.1.N120.dt5e-05.l2\endcsname{$4.71\times10^{-6}$}
\expandafter\def\csname num@cheb.orig.0.1.N40.dt0.002.l2\endcsname{$1.88\times10^{-4}$}
\expandafter\def\csname num@cheb.orig.0.1.N60.dt0.001.l2\endcsname{$9.41\times10^{-5}$}
\expandafter\def\csname num@cheb.orig.0.1.N80.dt0.0001.l2\endcsname{$9.41\times10^{-6}$}
\expandafter\def\csname num@cheb.orig.0.1.N80.dt0.0005.l2\endcsname{$4.70\times10^{-5}$}
\expandafter\def\csname num@ode.compat.0.01.l2\endcsname{$1.89\times10^{-3}$}
\expandafter\def\csname num@ode.compat.0.1.l2\endcsname{$1.47\times10^{-1}$}
\expandafter\def\csname num@ode.orig.0.01.l2\endcsname{$1.30\times10^{-3}$}
\expandafter\def\csname num@ode.orig.0.1.l2\endcsname{$1.02\times10^{-1}$}

\newcommand{\rc}[1]{\par\medskip\noindent\textbf{Comment.} \emph{#1}\par\smallskip\noindent\textbf{Response.} }
\setlength{\parskip}{4pt}
\setlength{\parindent}{0pt}

\begin{document}

\begin{center}
{\large\bfseries Response to the Reviewers}\\[4pt]
Manuscript symmetry-4596052 (resubmission)\\[2pt]
``A Chebyshev-Coefficient Physics-Informed Neural Network for the Allen--Cahn Equation: Exact Neumann Basis, A-Posteriori Energy Balance and Inverse Identification''\\
(first version: ``A Spectral--Physics-Informed Neural Network Hybrid with Energy-Constrained Training for the Allen--Cahn Equation: Forward and Inverse Problems'')
\end{center}

\medskip
Dear Editor, dear Reviewers,

We thank both reviewers for their careful reading. Their comments were correct, and following them up changed the paper substantially. Besides answering every point, we re-ran the one-dimensional experiments with a new evaluation protocol (the one-dimensional stress tests of the first version are superseded or removed, see the end of this letter) and audited every number of the first version against its raw data. We report the outcome honestly, including three conclusions of the first version that we now retract. Because most of the text was rewritten, the highlighted version of the manuscript shows in blue every paragraph, heading, table, caption and reference that is new or changed with respect to the first version; only the few passages carried over verbatim are black. The Lean~4 formalization of the first version has been removed from the manuscript; its role is taken by a short mathematical proof (Theorem~\ref{thm:balance}). All results of the training runs and classical baselines of this revision, in the revised text and in this letter, are generated by scripts from the archived raw results (repository given in the Data Availability Statement).

\section*{Summary of the revision}

\begin{enumerate}[leftmargin=*,itemsep=2pt]
\item {\sloppy \textbf{Initial condition.} The initial condition $x^2\cos(\pi x)$ violates the Neumann condition (Reviewer~1). All one-dimensional experiments of the main text now use the compatible condition $u_0=\frac12(\cos\pi x-\cos^2\pi x)$; Appendix~\ref{app:orig} repeats the B1 runs of the as-submitted network with the original one.\par}
\item \textbf{Certified reference and consistent norms.} The first-order finite-difference reference of the first version had wall rows that were only first-order accurate, and the Chebyshev row of the old Table~1 used a different norm from the PINN rows. All errors are now measured against a certified spectral reference (agreement to $\le10^{-8}$ under refinement and with an independent solver, Appendix~\ref{app:ref}) on one common grid.
\item \textbf{Revised network.} The Neumann condition is now imposed exactly through the basis $\varphi_k=T_k-\frac{k^2}{(k+2)^2}T_{k+2}$ (the classical Chebyshev--Galerkin Neumann basis, now cited), with Gauss--Lobatto collocation. The as-submitted variant is kept as a comparison on B1 and B2, in the mode-number, energy and joint-training inverse experiments, and in Appendix~\ref{app:orig}.
\item \textbf{Three benchmarks.} Because the original benchmark ($\varepsilon=0.01$, $T=1$) turned out to be reaction-dominated (a model without diffusion has an error of \N{ode.compat.0.01.l2}), we added a diffusion-active benchmark ($\varepsilon=0.1$) and an interface benchmark ($\varepsilon=0.05$, $T=5$).
\item \textbf{Controls and statistics.} Five seeds for the main configurations, a control run for every claimed effect, a bootstrap interval for the main ratio, stronger baselines (L-BFGS refinement), the classical solver at several resolutions, and the maximum bound and the energy for every network trained for the one-dimensional forward problem (the energy at three time resolutions for the Spectral-PINNs and on the grid of spacing $0.02$ for the vanilla PINNs).
\item \textbf{Energy balance.} A new a-posteriori energy balance for arbitrary trained trajectories, proved in Section~\ref{sec:balance} (Theorem~\ref{thm:balance}; the Lean~4 appendices of the first version are removed), with the residual and the boundary defect measured for every Spectral-PINN trained for the forward problem (Reviewer~1, point iii).
\item \textbf{Inverse problem.} A full ablation of the joint-training estimator, a new variable-projection estimator, a classical least-squares estimator on identical noise realizations, and an identifiability analysis.
\item \textbf{Retracted conclusions.} (a)~The large degradation of accuracy with $K$ (errors of order one) was caused by a boundary penalty whose effective weight was $K^{-4}$; a milder growth of the error with $K$ remains at a fixed training budget; a larger budget lowers the errors but does not remove this growth (Section~\ref{sec:results-k}). (b)~The ``robust $\sim38\%$ overestimation bias'' of the inverse estimate is an artifact of the starting value and the optimizer. (c)~The zero energy violations cannot be attributed to the energy penalty: the unpenalized network has none either. The novelty claims are toned down (Reviewer~1, point i), the title is changed accordingly, and the abstract is rewritten (under 200 words).
\item \textbf{Audit.} Of 77 numerical statements of the first version, 22 were wrong and 7 could not be verified from the stored data; all of them are corrected or withdrawn in the revision.
\end{enumerate}

\section*{Reviewer 1}

\rc{(i) The main idea of using neural networks to learn time-dependent spectral coefficients has already been developed in ``Spectrally adapted physics-informed neural networks for solving unbounded domain problems'' (2023). Moreover, penalizing energy increases during PINN training has also been explored in ``Energy dissipation preserving physics informed neural network for Allen--Cahn equations'' (2025). Therefore, neither spectral--PINN hybridization nor energy-dissipation penalization is novel in itself, and the manuscript does not clearly establish a substantive methodological advance from combining these existing ideas.}
We agree, and we thank the reviewer for both references, which we had missed. Both are now cited and discussed in the Introduction (Xia, B\"ottcher and Chou 2023 as the origin of time-only coefficient networks; K\"ut\"uk and Y\"ucel 2025 for the energy penalty), together with the related coefficient-space and structure-preserving networks (the spectral informed neural networks of Yu et al.\ 2026, neural spectral methods, neuro-spectral architectures, SAV-based networks, hard-constraint designs) and the comparisons of PINNs with classical solvers. We no longer claim either idea as new. The revised paper is positioned as a careful, reproducible evaluation, and its contributions are now stated as: (1)~the use of the classical, exactly Neumann-compatible Chebyshev--Galerkin basis of Shen (1995; now cited) as the output basis of the coefficient network, together with the finding that a soft Neumann penalty fails at large mode numbers in both of its scalings (it vanishes when normalized with $K$, and makes the loss ill-conditioned when it is not; Section~\ref{sec:results-k}); (2)~a forward study on three benchmarks with a certified reference and the classical solver as baseline; (3)~the a-posteriori energy balance of point~(iii); (4)~a controlled ablation of the energy penalty, which shows that it leaves the accuracy essentially unchanged in all regimes but one; (5)~the diagnosis of the inverse-problem ``bias'' and a comparison with a classical estimator; (6)~the demonstration of exact symmetry under $x\mapsto-x$ and $y\mapsto-y$ through the Fourier basis. The Discussion (Section~\ref{sec:niche}) now states plainly that the method is not a competitor of classical solvers on the forward problem.

\rc{(ii) The initial condition for the 1D problem does not satisfy the homogeneous Neumann boundary conditions, namely, $u_x(\pm1,0)=\mp2$.}
The reviewer is right, and we thank them for this remark, which triggered the complete re-examination. All one-dimensional experiments of the main text now use $u_0(x)=\frac12(\cos\pi x-\cos^2\pi x)=\sin^2(\pi x/2)\cos(\pi x)$, which satisfies $u_0'(\pm1)=0$ and has the same values as $x^2\cos(\pi x)$ at $x=0,\pm\frac12,\pm1$ (Section~\ref{sec:problem}). The certified reference eliminates the wall unknowns with the Neumann rows, so it is also correct for the incompatible condition (Appendix~\ref{app:ref}). Appendix~\ref{app:orig} repeats the B1 experiments of the as-submitted Spectral-PINN ($K=32$ and $64$, three seeds) with the original condition: with $K=32$ the errors are \N{orig.K32} (original) and \N{orig.K32c} (compatible) for the same seeds, and the conclusions that can be tested with the original condition do not depend on the choice (the exact Neumann basis of the revised network cannot represent it).

\rc{(iii) The Lean-verified results in Appendix A are not connected to a stability or convergence analysis of the proposed Spectral-PINN method. In particular, Theorem A1 assumes an exact gradient flow, but the manuscript does not establish that the trained neural-network approximation satisfies this assumption.}
This is correct: Theorem~A1 concerns exact gradient flows, and a trained network is not one. We have removed the Lean~4 formalization from the manuscript and replaced it by a short, self-contained mathematical proof that supplies the missing link (Theorem~\ref{thm:balance} in Section~\ref{sec:balance}; evaluation details in Appendix~\ref{app:balance}). For any trajectory $u=\sum_kc_k(t)T_k(x)$ with residual $r=u_t-\varepsilon^2u_{xx}+f(u)$ we prove the identity
\[
\frac{dE}{dt}=-\|u_t\|^2+\langle r,u_t\rangle+\varepsilon^2\bigl[u_xu_t\bigr]_{-1}^{1}\le\tfrac12\|r\|^2-\tfrac12\|u_t\|^2+\varepsilon^2\bigl[u_xu_t\bigr]_{-1}^{1},
\]
and its integrated form, so that the energy of \emph{any} trained network can grow only through its PDE residual and its boundary defect (the latter vanishes for the exact Neumann basis). The proof uses only that the energy is a polynomial in the coefficients, an integration by parts for polynomials and an elementary identity for the inner product. The residual and the defect are measured for every Spectral-PINN trained for the forward problem (Table~\ref{tab:balance} and Section~\ref{sec:results-balance}): over all of them the identity holds to \N{bal.id_err}, and for the revised $K=32$ networks the integrated residual term is \N{bal.int_r_hard} on average over the five seeds (at most \N{bal.int_r_hard_max}), against an actual energy decrease of \N{bal.dE_hard}. We state explicitly (Sections~\ref{sec:balance} and~\ref{sec:energy-discussion}) what this does not provide: it is an a-posteriori bound in terms of measured quantities, not a convergence proof, which would need a stability estimate for the PDE in terms of the optimization error.

\rc{Minor (i). The abstract is overly long and detailed, which obscures the main contribution.}
The abstract has been rewritten (under 200 words); it states the objective, the principal findings and the key limitations, and the experimental details are now in the main text.

\rc{Minor (ii). Energy decay alone does not guarantee physically admissible solutions ... I suggest adding a brief discussion of the maximum bound principle and whether it is respected in the numerical experiments.}
Thank you; the suggested reference (He and Pan, 2019) and the review of Du et al.\ (2021) are now cited, and $\max_{x,t}|u|$ is computed for every network trained on the one-dimensional forward problem (Table~\ref{tab:forward}, largest value per configuration, and Section~\ref{sec:results-balance}). The bound is not guaranteed: all revised Spectral-PINNs of B1 exceed it slightly (by at most \N{mbp.excess}, and by up to \N{mbp.excess_reg} in the poor-trajectory regimes), \N{mbp.n_van_over} of the \N{mbp.n_van} vanilla PINNs of B1 as well (by up to \N{mbp.van_excess}), and some inaccurate networks substantially ($\max|u|$ up to \N{kk.soft128.maxu}), although the energy of every one of them decreases monotonically. This illustrates the reviewer's point directly; Section~\ref{sec:results-balance} states that a guaranteed bound would need an additional mechanism such as a limiter, and the Discussion (Section~\ref{sec:energy-discussion}) that the maximum bound has to be checked separately.

\rc{Minor (iii). Please correct the range of $S(t)$ reported in line 390, as it is inconsistent with the values in Table 7.}
Corrected (Section~\ref{sec:2d-results}): at $t=1$ the violation lies between \N{s2d.smin} and \N{s2d.smax} in all configurations (RMS asymmetry at most \N{s2d.rms_max}, about three times the machine epsilon); the first version had written ``$10^{-31}$--$10^{-31.4}$''.

\section*{Reviewer 2}

\rc{1. The abstract states that the Spectral-PINN ``cuts the $L^2$ error by a factor of $\sim49$'' relative to a vanilla PINN, whereas Section 5 states ``a factor of about 47''. The authors must clarify the exact statistical basis of this factor and reconcile the abstract and conclusions.}
In the first version, 49 was the ratio of the five-seed mean errors and 47 the ratio for seed 42 alone; the two numbers should not have been mixed. In the revision the statistic behind every ratio of errors not explicitly attributed to individual seeds is the ratio of the seed means (Section~\ref{sec:protocol}), and the main ratio is given with a 95\% bootstrap interval. With the compatible initial condition and the revised network the ratio on B1 is \N{b1.ratio.van_fhK32} (interval \N{b1.ratio.van_fhK32.ci}, per-seed ratios \N{b1.ratio.van_fhK32.range}); with an additional L-BFGS stage for both networks it is \N{b1.ratio.lb}. On the diffusion-active benchmark B2 it is \N{b2.ratio.van_hard} with Adam alone, and after L-BFGS refinement the vanilla PINN is even slightly more accurate (\N{b2.lb_van.l2} against \N{b2.lb_hard_K32.l2}, one seed each) (Section~\ref{sec:results-forward}, Table~\ref{tab:forward}). The abstract and the conclusions now quote the same statistics and state these qualifications. We also stress that this is a ratio between two PINNs; the classical solver is more accurate than both.

\rc{2. Section 2.4 introduces the one-sided energy penalty with ``weight $\lambda_E$'', but no numerical value of $\lambda_E$ is given ... The authors should report $\lambda_E$ explicitly, and ideally include a small ablation showing the effect of $\lambda_E$ on both accuracy and energy violations.}
The weight was $\lambda_E=50$ in all penalized runs of the first version (10 checkpoints per unit time); it is now stated in Section~\ref{sec:energy-loss}. Table~\ref{tab:energy} gives the requested ablation, $\lambda_E\in\{0,1,10,50,500,5000\}$ with five seeds each: the mean error lies between \N{en.lam.min} and \N{en.lam.max} for all weights (seed-to-seed SD \N{en.lam.sd} at $\lambda_E=0$), no network has an energy increase on grids with spacing $0.02$, $0.002$ and $2\times10^{-4}$, and the penalty is nonzero in at most \N{en.lam.max_nonzero} training epoch per run. We also tested regimes with a poor trajectory (truncated training, a small learning rate, no warm-up, the as-submitted $K=64$ network). The penalty is nonzero in some training epochs of each of these regimes (and on B3), but it changes the result substantially only without warm-up, and there it makes it worse (\N{reg.nowarm.E50} against \N{reg.nowarm.E0}). Consequently, the claim of the first version that the penalty produced the zero violations is withdrawn: every one of the \N{en.n_all_spec} Spectral-PINNs trained for the forward problem in this study, penalized or not and accurate or not, has a monotone energy.

\rc{3. Figure 1 displays the discrete Ginzburg--Landau energy ... with and without the energy-constrained penalty, but the two curves are visually identical and completely overlapping. The manuscript does not explain what the reader should conclude from this.}
The curves overlap because the penalty is practically inactive (its term is positive in only \N{en.nonzero} training epochs of the five runs with $\lambda_E=50$): the unpenalized network already decreases its energy on every interval, by at least \N{en.min_step} per interval of length $0.02$, while the two trajectories differ by at most \N{en.max_diff} (Section~\ref{sec:results-energy}). Figure~\ref{fig:energy} now shows this difference explicitly (panel b) and the accuracy versus $\lambda_E$ (panel c), and the text states the conclusion: in this regime the penalty has no visible effect, and monotone energy is not evidence for it.

\rc{4. Table 1 shows that the classical Chebyshev spectral collocation method solves the problem in 0.16 s, whereas the Spectral-PINN takes 169 s and the vanilla PINN 919 s ... The authors should provide a much clearer statement of the method's intended niche.}
We agree with the reviewer's reading, and the revision says so explicitly. With consistent norms and a certified reference, Chebyshev collocation with a first-order time step is \N{b1.ratio.fh_K32_cheb80} times (B1) and \N{b2.ratio.hard_cheb80} times (B2) more accurate than the revised Spectral-PINN and about \N{b1.speed.fh_K32_cheb80} times faster (B1, timings on one machine); on B1 even a model without diffusion is within a factor \N{b1.ratio.ode_fh_K32} of the network, and on the interface benchmark B3 both PINNs fail while the classical solver does not. Section~\ref{sec:niche} now states what the experiments do and do not show: the method is not offered as a competitor on the forward problem; within the family of PINNs it is more accurate and cheaper than a space--time PINN trained with Adam (after L-BFGS refinement the accuracy advantage shrinks on B1 and vanishes on B2), structure that is a linear constraint on the coefficients can be imposed exactly, and the energy of the trained trajectory is controlled a posteriori. Possible niches (continuous-in-time differentiable representations, partially unknown physics, parametric families) are named as such, with the remark that we have not tested them and that our own inverse experiments give no support for the data-assimilation motivation in the simplest setting.

\rc{5. Section 3.2 identifies a robust $\sim38\%$ overestimation bias in recovering $\varepsilon$ and attributes it to ``loss weighting / warm-up initialization / observation design'', but no supporting ablation is provided.}
We carried out the ablation, and it showed that the ``bias'' is not a property of the estimator. With the joint-Adam procedure of the first version the estimate is $\hat\varepsilon\approx\N{inv.ratio_mean}\,\varepsilon_{\rm init}$ (range \N{inv.ratio_range} over \N{inv.n_sweep} runs) for every starting value from $0.5$ to $4\,\varepsilon_{\rm true}$, for two values of $\varepsilon_{\rm true}$ and for both networks; all 20 runs of the sweep of the first version that produced the ``38\% bias'' started from $2\varepsilon_{\rm true}$, and $0.7\times2=1.4$ is that bias. The size of the displacement is set by the learning rate of $\log\varepsilon$ (Table~\ref{tab:inverse2}); the observation design (10 to 160 points, 5 or 20 times), $K$ and the PDE weight change the estimate only within \N{inv.ab.design_range}$\,\varepsilon_{\rm true}$, and the estimate moves away from the true value when started there (Section~\ref{sec:results-inverse}, Tables~\ref{tab:inverse1} and~\ref{tab:inverse2}, Figure~\ref{fig:inverse}). The noise results of the first version ($40.2\%$, $40.2\%$, $40.0\%$) and the single run with $14.7\%$ came from a script whose settings differed from those stated and whose output had not been archived; they are withdrawn. As a remedy we introduce a variable-projection estimator (closed-form $\varepsilon^2$ for the current network); it recovers $\varepsilon$ from noise-free data to \N{inv.vp_err_005} and \N{inv.vp_err_01} at $\varepsilon_{\rm true}=0.05$ and $0.1$ (\N{inv.vp001.n0} at $\varepsilon_{\rm true}=0.01$), but it is biased under noise (\N{inv.vpn.e0.05.n0.01} for $\sigma=0.01$, $\varepsilon_{\rm true}=0.05$), whereas a classical least-squares estimator on the same noise realizations has \N{inv.cl5.e0.05.n0.01} and stays close to the Cram\'er--Rao bound (Table~\ref{tab:inverse3}). We also show that $\varepsilon=0.01$ is not identifiable from the data of this study at noise level $0.01$ (Cram\'er--Rao standard error \N{ident.se001}), so that the noise experiments at $\varepsilon=0.01$ of the first version ($\sigma=0.01$ and $0.05$) were not meaningful tests.

\section*{Further corrections not requested by the reviewers}

\begin{itemize}[leftmargin=*,itemsep=2pt]
\item The reference solution of the first version was first-order accurate at the walls; it is replaced by a certified spectral reference (Appendix~\ref{app:ref}). The Chebyshev row of the old Table~1 used a different norm from the PINN rows; all rows now use the same grid and norm.
\item The statements on the number of threads and on the hardware were inaccurate; Section~\ref{sec:protocol} now describes the actual setup, and the times in Table~\ref{tab:forward} come from separate timing runs on one machine.
\item The ``DoF'' column of the old Table~1 mixed parameter counts and coefficient counts and has been removed.
\item The statement that more Chebyshev nodes are ``unambiguously beneficial'' in the classical solver was not supported by the stored numbers; the solver is now reported with consistent norms (at $N=80$ with two time steps on B1, and at two resolutions on B3), and on B1 and B2 its error is dominated by the time step.
\item The long-horizon and large-$K$ stress tests of the first version were run with the as-submitted boundary penalty and are superseded by Section~\ref{sec:results-k} and benchmark B3.
\item The smaller-$\varepsilon$ stress test of the first version (old Table~4, $\varepsilon$ down to $0.001$) has been removed: its Chebyshev column used an inconsistent norm, the Chebyshev resolution was not scaled with $1/\varepsilon$ as stated, and the fixed 200-point evaluation grid does not resolve the interface at the smallest $\varepsilon$.
\end{itemize}

We hope that the revised manuscript meets the reviewers' concerns, and we thank them again for remarks that substantially improved the work.

\medskip
Sincerely,\\
Zhanars Abdiramanov, on behalf of all authors

\end{document}
